

















































































































































\section{Introduction}
We study a class of optimization problems ubiquitous in theoretical computer science, machine learning,  quantum information theory and statistics. The programs we consider are continuous optimization problems over matrix groups. More precisely, they can be posed as Euclidean norm minimization over the closure of a group orbit. The programs span two historically distinct contexts: In one context, the optimization problems are convex, and in the other they are not convex but rather \emph{geodesically convex} on a suitable manifold.


The \emph{commutative setting}, in which the underlying group is Abelian, captures matrix scaling, matrix balancing and array scaling, which arise in scientific computing and optimal transport \cite{cuturi2013sinkhorn, parlett1971balancing}. Such problems fall into the framework of unconstrained geometric programming. Though these problems are convex, there are at least two reasons to study them further. Firstly, they are of such practical importance that speed matters. Na\"ively applying powerful algorithms like ellipsoid and interior point methods can be impractically slow. Hence, it is important to understand when faster methods can succeed. Matrix scaling and balancing, in particular, have enjoyed some success stories~- there are fast algorithms to obtain high precision solutions \cite{cohen2017matrix,allen2017much}, and there are more general upper bounds \cite{burgisser2020interior}. Secondly, the algorithms developed for the commutative setting are candidates for generalization to our second setting, which takes place in the less well-understood arena of geodesically convex optimization.


The second context, which we call the \emph{noncommutative setting}, arises when the underlying group is non-Abelian. The noncommutative setting captures problems like operator and tensor scaling \cite{garg2016deterministic,burgisser2017alternating},  the quantum marginal problem \cite{burgisser2018efficient} and statistical estimators such as Tyler's M estimator \cite{franks2020rigorous} and maximum likelihood estimates for matrix and tensor normal models \cite{amendola2020invariant}. Deciding whether the value of the optimization problem is zero or not is equivalent to deciding a central polynomial identity testing (P.I.T.) problem in invariant theory known as the \emph{null cone} problem. It is hoped that efficient optimization algorithms will result in efficient algorithms for the null-cone problem.
One approach to complexity lower bounds, geometric complexity theory, suggests that these P.I.T. problems should be in $\P$ \cite{mulmuley2017geometric, garg2019search}, and the optimization approach has resulted in polynomial time algorithms in some cases \cite{garg2016deterministic, allen2018operator}. The optimization problems that arise in the noncommutative setting are not convex in the Euclidean sense, but rather \emph{geodesically convex}, a notion of convexity on a Riemannian manifold. Currently, the only implementable algorithms for geodesically convex optimization are analogues of gradient descent and trust region methods \cite{absil2008optimization, zhang2016first,allen2018operator}. There are, as of yet, no efficiently implementable geodesically convex counterparts to the interior point or cutting plane methods.


In both the commutative and noncommutative settings, algorithms are typically analysed using two quantities. One is \emph{diameter}, or how far approximate minimizers can be from the origin. The other is a geometric measure of well-conditionedness known as \emph{margin} (or \emph{gap} in the noncommutative case), which has several variants in the literature and appears in two primary ways. Firstly, the smaller the margin, the higher the degree of precision required to decide if the value of the optimization problem is zero or not \cite{gradflow, gurvits2004combinatorial}. Secondly, the larger the margin, the smaller the diameter \cite{singh2014entropy,straszak2017computing,gradflow,burgisser2020interior}. In this paper we show the following:


\begin{enumerate} 
\item[i)] In the commutative setting, and in particular for array scaling, approximate minimizers for the functions we study can have doubly exponential condition number. That is, the problems have exponential diameter. As a consequence, popular classes of algorithms such as gradient descent and trust region methods cannot produce high-precision solutions in polynomial time in general. This result applies in the noncommutative setting as well, which provides evidence that even cutting plane methods are unlikely to produce high-precision solutions in polynomial time. This shows it is necessary to develop powerful methods like the interior point method in the geodesically convex setting. 
\item[ii)] In the commutative and noncommutative settings, we study the margin and gap, respectively, which appear in running time bounds for all existing algorithms. We prove that these measures can be exponentially small in the input size for several problems including array scaling and tensor scaling. In the commutative case, this gives evidence that existing algorithms for array scaling do not run in near-linear time. In the noncommutative case, our results show that margin-based analyses like \cite{gradflow} cannot prove polynomial time guarantees for deciding the null cone problem for tensor scaling using trust region methods. 
\end{enumerate} 


We use the remainder of the introduction to describe both settings in more detail, state our main results precisely, and discuss previous work. For both the commutative and noncommutative settings, we proceed in the following order. We start with an introduction and motivation of the setting, continue with diameter bounds and afterwards treat bounds on the margin and gap, respectively. We end each setting with a short discussion of the main proof techniques.



\subsection{The commutative setting: matrix scaling and its relatives}


\paragraph{Matrix scaling and array scaling.}

Consider the matrix scaling problem: given a nonnegative matrix $A$, find nonnegative diagonal matrices $X, Y$ such that $XA Y$ is doubly stochastic (i.e. has row and column sums equal to one). The matrices, if they exist, can be found by the exceedingly simple and fast alternating minimization method known as \emph{Sinkhorn's algorithm}. It is frequently used in practice, e.g. for quickly approximating the solution to optimal transport problems \cite{cuturi2013sinkhorn}.

Like all other algorithms for matrix scaling,  Sinkhorn's algorithm is typically analyzed through optimization. One finds that $X$ and $Y$ are $e^{\diag(x)}, e^{\diag(y)}$, where $x,y \in \RR^n$ are solutions to the following optimization problem: 
\begin{align}\inf_{x, y \in \RR^n} \sum A_{ij}e^{x_i + y_j - \bar{x} - \bar{y}}\label{eq:matrix-capacity}\end{align}
for $\bar{z}:=\frac{1}{n} \sum z_i$ (c.f. \cite{kalantari1996complexity}). Moreover, the infimum is greater than zero if and only if $A$ is \emph{approximately scalable}, i.e. the row and column sums of $XAY$ can be made arbitrarily close to one for $X, Y$ nonnegative, diagonal.


More generally, given a finite set $\Omega \subseteq \RR^m$ and a nonnegative function $p:\Omega \to \RR_{\geq 0}$, define the \emph{capacity} \cite{gurvits2004combinatorial} as the value of the unconstrained geometric program
\begin{gather}
\capa(p):= \inf_{x \in \RR^m} f_p(x) := \inf_{x \in \RR^m} \sum_{\omega \in \Omega} p_\omega e^{\omega \cdot x}.\label{eq:abelian}
\end{gather}
The capacity is positive if and only if zero is in the \emph{Newton polytope} $\conv(\supp p)$. Matrix scaling arises when $m = 2n$ and $\Omega = \{(\eps_i, \eps_j): i,j \in [n]\}$ for $\eps_k:=e_k - \frac{1}{n} \id_n$, where $e_k \in \RR^n$ is the $k^{th}$ canonical unit vector and $\id_n \in \RR^n$ denotes the all-ones vector. In this case \cref{eq:abelian} reduces to precisely \cref{eq:matrix-capacity}, and $\|\nabla \log f_p (x)\|$ measures the deviation of $p$ from doubly stochastic.

Matrix balancing, in which we instead wish to find a scaling for which the $i^{th}$ row and column sum match, arises when $m = n$ and $\Omega = \{e_i - e_j: i\neq j \in [n]\}$. When $m = 3n$ and $\Omega = \{(\eps_i, \eps_j, \eps_k): i,j,k \in [n]\}$ we obtain the $3$-dimensional \emph{array scaling problem}. In analogy to matrix scaling, in array scaling one has an array $p$ of numbers in $(\RR^n_{\geq 0})^{\ot 3}$ and seeks positive vectors $X,Y,Z \in \RR^n_{\geq 0}$ so that the array $q$ with entries $q_{ijk} =p_{ijk}X_i Y_j Z_k$ is \emph{tristochastic}. That is, the sum over every \emph{slice} is equal to one, i.e.
    $\sum_{j,k} q_{i_0,j,k} = \sum_{i,k} q_{i,j_0,k} = \sum_{i,j} q_{i,j,k_0} = 1$
for all $i_0,j_0,k_0 \in [n]$. If it is possible to satisfy these equations to arbitrary precision we say $p$ is approximately scalable. As for matrix scaling, $p$ is approximately scalable if and only if $\capa(p) > 0$. In the same manner, we obtain $d$-dimensional array scaling for $m=dn$ and
	\begin{equation}\label{eq:defnOmega-n-d}
	\Omega = \Omega_{n,d} := \big\lbrace \eps_i \colon i \in [n] \big\rbrace^d \subseteq \big( \RR^n \big)^d.
	\end{equation}
We can think of subsets of $\Omega_{n,d}$ as $d$-uniform, $d$-partite hypergraphs. Up to an additive shift by $- \frac{1}{n} \id_{nd}$, the elements of $\Omega_{n,d}$ are indicator vectors of the edges in such hypergraphs. For $d = 2$, the matrix $p$ is scalable if and only if the bipartite graph corresponding to $\supp p$ contains a perfect matching, but this is not the case for $d \geq 3$ (indeed, $d$-partite hypergraph matching is $\NP$-hard).



\paragraph{Algorithms for array scaling.}

Array scaling serves the same role for speeding up multimarginal transport as matrix scaling for optimal transport, and yet again there is a simple and fast alternating minimization algorithm that produces $\eps$-tristochastic scalings in time $O(1/\eps^2)$ \cite{altschuler2020polynomial,lin2019complexity}.
Moreover, algorithms to approximate the capacity arise in varied settings including radial isotropic position \cite{hardt2013algorithms}, entropy maximization \cite{straszak2017computing}, and approximate counting \cite{anari2018log}.  
 

It is natural to ask if there are \emph{high-precision} algorithms for array scaling with $\log(1/\eps)$ dependence on the error and linear or mild dependence on the number of nonzero entries. For matrix scaling and matrix balancing, several works have shown that trust regions and interior point methods \emph{can} obtain such guarantees \cite{cohen2017matrix,allen2017much}.
Our work is concerned with whether the performance of such algorithms carries over to array scaling and the computation of the capacity in general.



\subsubsection{Diameter lower bounds}

Guarantees for many iterative algorithms in convex optimization require \emph{diameter bounds}, or bounds on the distance $R$ from the starting point to an $\eps$-approximate solution. Trust region methods, also called \emph{box-constrained Newton's method}, are iterative algorithms that, at each step, move to the best solution within a typically small distance $D$ of the previous solution. By their nature, trust region methods take at least $R/D$ steps to produce an $\eps$-approximate solution. Gradient descent for Lipschitz functions also depends quadratically on a diameter bound, and cutting plane methods typically use diameter bounds to control the volume of a starting region.

\paragraph{Known diameter upper and lower bounds.}
For matrix scaling and matrix balancing, it has been shown in \cite{cohen2017matrix} that one may take $R = O(n \log(w_A/\eps))$, where $w_A$ is the ratio between the sum of the entries of the matrix and the least nonzero entry. For $3$-dimensional array scaling, the best upper bound of which we are aware is $R = O(n^{3/2} 2^{6n} \log (1/\eps)),$ which follows from the general upper bound of \cite{straszak2017computing} on diameter bounds for unconstrained geometric programming. There is also a diameter bound for array scaling in the multimarginal transport context that is polynomial in the input size assuming the tensor has no nonzero entries \cite{lin2019complexity}.

Regarding diameter \emph{lower} bounds, in the context of computing maximum entropy distributions it was shown that there is some bounded set $\Omega \subset \ZZ^m$ in a $\poly(m)$ size ball such that there are \emph{no} $\eps$-approximate minimizers of norm $\poly(m, \log 1/\eps)$ for $f_p$ as in \cref{eq:abelian}  \cite{straszak2017computing}. 

\paragraph{Main theorem.}
Where do the polynomial diameter bounds for matrix scaling (i.e. $2$-dimensional array scaling) transition to the superpolynomial diameter bounds for general $\Omega$? We show that this transition takes place in the next simplest problem, the $3$-dimensional array scaling problem.

\begin{thm}\label{thm:diameter}
There is an absolute constant $C > 0$ and an array $p_{ijk} \in (\RR_{\geq 0}^n)^{\ot 3}$ with $O(n)$ nonzero entries, each of bit-complexity $O(n)$, that satisfies the following property. For all $0 <\eps \leq  \exp(- C n^2 \log n)$ and $(x,y,z) \in \RR^{3n}$, if
	$$f_p (x,y,z) \leq \capa(p) + \eps $$
	then $\norm{(x,y,z)}_2 = \Omega\left(2^{n/3}\log(1/\eps)\right).$
\end{thm}

To emphasize that the difficulties do not lie in an additive vs multiplicative approximation, we remark that our array $p$ has unit sum and $\capa(p) = 1/2$. By a simple duplication trick, the same bound holds for $d$-dimensional array scaling with $d \geq 3$; see \cref{cor:diameter-d}.



\paragraph{Implications of \cref{thm:diameter} and relation to the literature.}

Theorem~\ref{thm:diameter} shows that trust region methods for array scaling with polynomial step size cannot provide high-precision solutions in $\poly(n,\log(1/\eps))$ time for
$d \geq 3$. Moreover, gradient descent on the Lipschitz convex function $\log f_p$ has a bounded step size, and so also cannot provide high precision solutions in polynomial time.

In \cite[Section 2.1]{straszak2017computing} the authors ask whether there is $\Omega$ whose elements are Boolean (up to an additive shift) with a superpolynomial diameter lower bound. As subsets of $\Omega_{n,d}$ are automatically of this form, we answer their open problem in the affirmative. Our lower bound on $\log R$ is tight up to constant factors by the diameter upper bound from \cite{straszak2017computing} mentioned above; moreover the logarithmic dependence on $\eps$ is best possible. Determining the correct constant in the exponent is an interesting open direction. We believe that that the requirement that $\eps$ is very small is an artifact of our specific construction and proof strategy, and thus can probably be relaxed significantly.


































































































































































































































































































































































































































































































































































































































































\setcounter{section}{2}\section{Diameter bounds in the commutative case}\label{sec:diameterComm}


In this section we describe an array such that all approximate scalings are very ill conditioned, proving \cref{thm:diameter}. Let us define the diameter bound.
\begin{dfn}\label{dfn:diameterBound} Let $\eps \to 0$ and $f:\RR^m \to \RR$. The \emph{diameter bound} $D_f(\eps)$ is defined as the infimum over $R > 0$ such that 
$$ \inf_{\|x\| \leq R} f(x) \leq \eps + \inf_{x \in \RR^m}f(x).$$
\end{dfn}
Thus, \cref{thm:diameter} is equivalent to the statement that $D_f(\eps) = \Omega(2^{n/3} \log (1/\eps)$ for $\eps \leq e^{- C n^2 \log n}$. We now give a proof outline for \cref{thm:diameter}. 

\subsection{Proof outline}\label{subsec:outline}

The high-level intuition applies not only to array scaling but to the capacity in general. Recall that the array scaling capacity is 
$$ \inf_{x \in \RR^{3n}}  \sum_{\omega \in \Omega} p_{\omega} e^{\omega \cdot x} $$ for $\Omega =  \Omega_{n,3} = \{e_i - \frac1n \id_n: i \in [n]\} \subseteq \RR^{3n}$. We build both the support $\Omega'\subseteq \Omega_{n,3}$ and the entries $p$ in the following way. We construct a set $\Omega_0 \subseteq \Omega_{n,3}$, another element $\omega \in \Omega_{n,3}$, and an array $q$ with the following properties.
\begin{enumerate}
\item The set $\Omega_0 \subseteq \Omega_{n,3}$ should be the support of a tristochastic array $q$. 
\item The affine hull of $\Omega_0$, should have codimension one\footnote{This will not quite apply in our setting, because $\aff(\Omega_{n,3})$ is not full-dimensional. Instead, $\aff (\Omega_0)$ will be codimension one in $\aff(\Omega_{n,3})$.} in $\RR^{3n}.$ 
\item The origin is in the relative interior of $\conv(\Omega_0)$. Note that the origin is already in $\conv(\Omega_0)$ by the tristochasticity of $q$.
\item The vector $\omega \in \Omega_{n,3}$ should be at a very small, but positive, distance $\eta$ from $\aff (\Omega_0)$. Note that this already implies that the facet gap of $\Omega_0 \cup \omega$ is small.
\end{enumerate}
Finally, we define the entries of $p$ by $p|_{\Omega_0} = \frac{1}{2} q$, $p_\omega = \frac{1}{2}$, and $p_\omega = 0$ elsewhere. Assuming we have found $p$ according to this process, we now give intuition for the diameter bound.

Let $v$ be the projection of $\omega$ to the orthogonal complement of $\aff(\Omega_0)$. Intuitively, the capacity is only approximately attained by vectors very far in the $-v$ direction. Indeed, first note that $\capa(p) = 1/2$, because $\capa(q) = 1$ by tristochasticity, $\capa(p) \geq \frac{1}{2} \capa(q) = \frac{1}{2}$, and $f_p(- tv/\|v\|) = \frac{1}{2} + e^{ - \eta t}$ so $f_p(- tv/\|v\|)$ tends to $\frac 12$. However, $f_p( - t v/\|v\|)$ tends to $\frac 12$ slowly if $\eta$ is small. Indeed, $f_p( - t v/\|v\|) \leq \frac{1}{2}(1 + \eps)$ only if $t \geq \frac{1}{\eta} \log (1/\eps)$. 

To conclude rigorously that the capacity is only approached by vectors very far in the $-v$ direction, we must rule out directions with nonzero components in $\aff(\Omega_0)$. For this, we must use the assumption that $0$ is rather deep in the relative interior of $\conv(\Omega_0)$. If this is the case, then any $\eps$-approximate minimizer must have a bounded component in $\aff(\Omega_0)$, for otherwise the contribution to $f_p$ from the elements of $\Omega_0$ alone will be larger than $\frac{1}{2} + \eps$. 

The remainder of the section will be concerned with the construction of a subset $\Omega_0$, an array $q$, and  an element $\omega$ with these properties.


\subsection{The construction}\label{subsec:diameter-constr}

We construct the subset $\Omega_0$ from a directed graph $D$ on $[n]$, which we will determine later. If $i,j$ is an edge in $D$, then $\Omega_0$ includes the elements $(\eps_i, \eps_i, \eps_j)$ as well as the three cyclic permutations of it. That is, 
$$ \Omega_0 = \{(\eps_j, \eps_i, \eps_i), (\eps_i, \eps_j, \eps_i), (\eps_i, \eps_i, \eps_j): ij \in E(D)\}.$$
We now describe the graph, as seen in \cref{fig:graph}.

\begin{figure}
\centering
\begin{tikzpicture}
%
\GraphInit[vstyle=Classic]
%
\useasboundingbox (0,0) rectangle (15.0cm,15.0cm);
%
\Vertex[L=\hbox{$\overline{w}_{l-1}$},x=7.468cm,y=1.6336cm]{v0}
\Vertex[L=\hbox{$\overline{w}_{l}$},x=8.1205cm,y=0.0cm]{v1}
\Vertex[L=\hbox{$r$},x=7.3452cm,y=11.4526cm]{v2}
\Vertex[L=\hbox{$u_1$},x=9.6222cm,y=12.531cm]{v3}
\Vertex[L=\hbox{$u_{l-1}$},x=13.5892cm,y=14.2394cm]{v4}
\Vertex[L=\hbox{$u_{l-2}$},x=11.7656cm,y=13.4716cm]{v5}
\Vertex[L=\hbox{$u_{l}$},x=15.0cm,y=14.8352cm]{v6}
\Vertex[L=\hbox{$v_1$},x=5.2927cm,y=12.7084cm]{v7}
\Vertex[L=\hbox{$v_{l-1}$},x=1.4984cm,y=14.4533cm]{v8}
\Vertex[L=\hbox{$v_{l-2}$},x=3.3291cm,y=13.6941cm]{v9}
\Vertex[L=\hbox{$v_{l}$},x=0.0cm,y=15.0cm]{v10}
\Vertex[L=\hbox{$w_1$},x=7.1865cm,y=8.8119cm]{v11}
\Vertex[L=\hbox{$w_{l-1}$},x=5.1617cm,y=2.089cm]{v12}
\Vertex[L=\hbox{$w_{l-2}$},x=6.5507cm,y=3.5411cm]{v13}
\Vertex[L=\hbox{$w_{l-3}$},x=6.9099cm,y=6.2032cm]{v14}
\Vertex[L=\hbox{$w_{l}$},x=4.1288cm,y=0.8072cm]{v15}
%
\Edge[label=\hbox{$2 + 1$}, labelstyle={sloped,below,fill=\thepagecolornone}](v0)(v1)
\Edge[label=\hbox{$2 - \frac{1}{2}$}, labelstyle={sloped,below,fill=\thepagecolornone}](v0)(v13)
\Edge[label=\hbox{$2 + (-\frac{1}{2})^{l-1}$}, labelstyle={scale=0.8,sloped,above,fill=\thepagecolornone}](v2)(v3)
\Edge[label=\hbox{$2 + (-\frac{1}{2})^{l-1}$}, labelstyle={scale=0.8,sloped,below,fill=\thepagecolornone}](v2)(v7)
\Edge[label=\hbox{$2 + (-\frac{1}{2})^{l-2}$}, labelstyle={scale=0.8,sloped,above,fill=\thepagecolornone}](v2)(v11)
\Edge[style=dotted](v3)(v5)
\Edge[label=\hbox{$2 - \frac{1}{2}$}, labelstyle={sloped,above,fill=\thepagecolornone}](v4)(v5)
\Edge[label=\hbox{$2 + 1$}, labelstyle={sloped,above,fill=\thepagecolornone}](v4)(v6)
\Edge[style=dotted](v7)(v9)
\Edge[label=\hbox{$2 - \frac{1}{2}$}, labelstyle={sloped,below,fill=\thepagecolornone}](v8)(v9)
\Edge[label=\hbox{$2 + 1$}, labelstyle={sloped,below,fill=\thepagecolornone}](v8)(v10)
\Edge[style=dotted](v11)(v14)
\Edge[label=\hbox{$2 - \frac{1}{2}$}, labelstyle={sloped,above,fill=\thepagecolornone}](v12)(v13)
\Edge[label=\hbox{$2 + 1$}, labelstyle={sloped,above,fill=\thepagecolornone}](v12)(v15)
\Edge[label=\hbox{$2 -\frac{1}{2}$}, labelstyle={sloped,above,fill=\thepagecolornone}](v13)(v14)
%
\end{tikzpicture}

\caption{The graph $D_l$ from \cref{dfn:graph} with the edge labels proportional to the edge labeling $q$ in \cref{it:stoch} of \cref{lem:diameter} (the constant factor $1/6n$ is omitted for readability). We have also omitted the directions, which are all towards the root $r$.
}\label{fig:graph}
\end{figure}




\begin{figure}
\centering
$\left[\begin{array}{c|c|c|c|c}
\scalebox{1.5}0 &
\begin{array}{ccccc}
\cellcolor{blue!30}\mathrlap{A}{\phantom{\ddots}} & \cellcolor{green!30} I & \hdots & 0 \\
\mathrlap{0}{\phantom{\ddots}} & \cellcolor{blue!30}\ddots & \cellcolor{green!30}\ddots & \smash{\vdots} \\
\mathrlap{\vdots}{\phantom{\ddots}} & \ddots & \cellcolor{blue!30}\ddots & \cellcolor{green!30}I \\
\mathrlap{0}{\phantom{\ddots}} & \smash{\hdots} & \mathrlap{0}{\phantom{\ddots}} & \cellcolor{blue!30}\mathrlap{A}{\phantom{\ddots}} 
\end{array}
& \scalebox{1.5}0 & \scalebox{1.5}0 & 
\begin{array}{c}
\mathrlap{0}{\phantom{\ddots}} \\
\mathrlap{\vdots}{\phantom{\ddots}}  \\
\mathrlap{\vdots}{\phantom{\ddots}}  \\
\cellcolor{green!30}\mathrlap{I}{\phantom{\ddots}} 
\end{array}
\\
\hline
\scalebox{1.5}0 & \scalebox{1.5}0 & 
\begin{array}{ccccc}
\cellcolor{blue!30}\mathrlap{A}{\phantom{\ddots}} & \cellcolor{green!30} I & \hdots & 0 \\
\mathrlap{0}{\phantom{\ddots}} & \cellcolor{blue!30}\ddots & \cellcolor{green!30}\ddots & \smash{\vdots} \\
\mathrlap{\vdots}{\phantom{\ddots}} & \ddots & \cellcolor{blue!30}\ddots & \cellcolor{green!30}I \\
\mathrlap{0}{\phantom{\ddots}} & \smash{\hdots} & \mathrlap{0}{\phantom{\ddots}} & \cellcolor{blue!30}\mathrlap{A}{\phantom{\ddots}} 
\end{array}
& \scalebox{1.5}0 & 
\begin{array}{c}
\mathrlap{0}{\phantom{\ddots}} \\
\mathrlap{\vdots}{\phantom{\ddots}}  \\
\mathrlap{\vdots}{\phantom{\ddots}}  \\
\cellcolor{green!30}\mathrlap{I}{\phantom{\ddots}} 
\end{array}\\
\hline
\scalebox{1.5}0 & \scalebox{1.5}0 & \scalebox{1.5}0 & 
\begin{array}{ccccc}
\cellcolor{blue!30}\mathrlap{A}{\phantom{\ddots}} & \cellcolor{green!30} I & \hdots & 0 \\
\mathrlap{0}{\phantom{\ddots}} & \cellcolor{blue!30}\ddots & \cellcolor{green!30}\ddots & \smash{\vdots} \\
\mathrlap{\vdots}{\phantom{\ddots}} & \ddots & \cellcolor{blue!30}\ddots & \cellcolor{green!30}I \\
\mathrlap{0}{\phantom{\ddots}} & \smash{\hdots} & \mathrlap{0}{\phantom{\ddots}} & \cellcolor{blue!30}\mathrlap{A}{\phantom{\ddots}} 
\end{array} & 
\begin{array}{c}
\mathrlap{0}{\phantom{\ddots}} \\
\mathrlap{\vdots}{\phantom{\ddots}}  \\
\mathrlap{\vdots}{\phantom{\ddots}}  \\
\cellcolor{green!30}\mathrlap{I}{\phantom{\ddots}} 
\end{array}\\
\hline
\begin{array}{cc} \cellcolor{blue!30}\mathrlap{A}{\phantom{\ddots}} & \cellcolor{green!30}\mathrlap{I}{\phantom{\ddots}}  \\
0 & \cellcolor{blue!30}\mathrlap{A}{\phantom{\ddots}} 
\end{array} & \scalebox{1.5}0 & 
\scalebox{1.5}0
& \begin{array}{cccc} \mathrlap{0}{\phantom{\ddots}} & \mathrlap{0}{\phantom{\ddots}} & \mathrlap{0}{\phantom{\ddots}} & \hdots \\
\mathrlap{0}{\phantom{\ddots}} & \mathrlap{0}{\phantom{\ddots}} &  \cellcolor{green!30}\mathrlap{I}{\phantom{\ddots}} & \hdots
\end{array} 
& \begin{array}{c}\mathrlap{0}{\phantom{\ddots}}\\ \mathrlap{0}{\phantom{\ddots}}\end{array}
\end{array}\right]
$
\caption{The matrix $M$ written in the reordered basis described before \cref{lem:diameter}. From the left, the five groups of columns correspond to the $\overline{w}'s$, the $u's$, the $v's$, the $w's$, and $r$ among the vertices of $D_l$. As such the dimensions of the five column groups, from left, are $3 \cdot 2, 3  (l-1), 3  (l-1), 3  (l-1), 3$, and the dimensions of the four groups of rows from top are $3  (l-1), 3  (l-1), 3  (l-1), 3\cdot 2$. $A$ is as in \cref{eq:a-matrix} and $I$ is the $3\times 3$ identity matrix.}\label{fig:matrix}
\end{figure}


\begin{figure}
\centering
\begin{tikzpicture}[
            > = stealth, 
            shorten > = 1pt, 
            auto,
            node distance = 3cm, 
            semithick 
        ]
\GraphInit[vstyle=Classic]
%
\useasboundingbox (0,-1.5) rectangle (5cm,1.5cm);
%
\Vertex[L=\hbox{$\;$},x=2.5cm,y=0cm]{v0}
\node at (2.5cm,.5cm) [draw,draw=none] {$v$};
\Vertex[L=\hbox{$\;$},x=5.0cm,y=0cm]{v1}
\Vertex[L=\hbox{$\;$},x=0.0cm,y=0.0cm]{v2}
%
\path[->] (v2) edge node {$q_1$} (v0);
\path[->] (v0) edge node {$q_2$} (v1);

\end{tikzpicture}

\caption{If $v$ is a vertex of $D_l$ with edges weighted $q_1$ and $q_2$ incident to it, then the column $v,i$ of $M$ for $i \in [3]$ sums to $q_1 + 2q_2$. That is, the incoming edge contributes its weight and the outgoing edge contributes twice its weight.}
\label{fig:column-sum}
\end{figure}

\begin{dfn}\label{dfn:graph}
The graph $D_l = (W,E)$ is a directed tree with $l + 1$ levels, where the root is on the $0^{th}$ level and the leaves are on the $l^{th}$ level. The tree is constructed as follows.
\begin{itemize}
\item All the edges are directed towards the root and are between adjacent levels. 
\item The root has three children, and on the $l-1$ levels below the root every node has one child. 
\item Additionally, one of the vertices on level $l-2$ has an additional child which has its own child. 
\end{itemize}
Explicitly, the vertices $W$ and edges $E$ are given by
\begin{align*}W &= \{u_i, v_i, w_i: i \in [l]\} \cup \{w_0:=u_0:=v_0:=r, \bar{w}_{l-1}, \bar{w}_l\}.\\
 E &= \{u_{i}u_{i-1}, v_{i}v_{i-1}, w_{i}w_{i-1}: i \in [l] \} \cup \{\bar{w}_{l-1} w_{l-2}, \bar{w}_l \bar{w}_{l-1}\}.\end{align*}
\end{dfn}
Note that $D_l$ has $3(l + 1)$ vertices so we set $n = 3(l + 1)$. Thus $D_l$ has $3 l + 2$ edges and so $|\Omega_0| = 3(3l + 2) = 3n - 3.$ It is helpful to construct the matrix $M$ whose set of rows is $\Omega_0$. To make the matrix sparser, first replace $\eps_i$ by $e_i$ by restricting the minimization to the subspace $\sum x_i = \sum y_i = \sum z_i = 0$, which is without loss of generality. We define $\Omega_0' \subseteq \RR^{3n}$ to be $\Omega_0$ but with each $(\eps_i, \eps_j, \eps_k)$ replaced by $(e_i, e_j, e_k)$; define $\Omega'_{n,3}$ similarly and define $p_{(e_i, e_j, e_k)}:=p_{(\eps_i, \eps_j, \eps_k)}$. Then
$$ \inf_{x \in \RR^{3n}} \sum_{\omega \in \Omega_{n,3}} p_{\omega} e^{(\eps_i, \eps_j, \eps_k) \cdot x} =  \inf_{\underset{\sum x_i = \sum y_i = \sum z_i = 0}{x,y,z \in \RR^{n}}} \;\sum_{\omega \in \Omega'_{n,3}} p_{\omega} e^{(e_i, e_j, e_k) \cdot (x,y,z)}.$$
 Moreover, when we write the matrix $M$, it is easier to write the vector $(x,y,z)$ in the order $(x_1,y_1,z_1, x_2,y_2,z_2, \dots )$ instead of the order $(x_1, \dots, x_n, y_1, \dots, y_n, z_1, \dots, z_n)$. With this ordering, the matrix $M$ with rows in $\Omega_0'$ is a block matrix $M$ with blocks of size $3$, with $n-1$ block rows, and with $n$ block columns. Each block row corresponds to an edge in the directed graph $D_l = (W, E)$ on $n = 3(l+1)$ vertices. If $e \in E$ is an edge from $i \to j$, then the $e^{th}$ row of $M$ has the matrix 
\begin{gather}A = \begin{bmatrix} 0 & 1 & 1 \\ 1 & 0 & 1 \\ 1 & 1 & 0 \end{bmatrix} \label{eq:a-matrix}\end{gather}
 in the $i^{th}$ block entry and 
 $$ I = \begin{bmatrix} 1 & 0 & 0 \\ 0 & 1 & 0 \\ 0 & 0 & 1 \end{bmatrix} $$ in the $j^{th}$ block entry and zeroes elsewhere. See \cref{fig:matrix} for a portrayal of the whole matrix $M$.
 
 
The first three properties for $\Omega_0$ in the proof plan translate to the following three claims about $M$. The first relates to the tristochasticity of $q$, the second to the codimension of $\aff(\Omega_0')$ in the subspace $(\id_n^\perp)^3$, and the third to the depth of the point $\frac{1}{n}\id_{3n}$ in $\conv(\Omega_0')$. 
 
 
 \begin{lem}\label{lem:diameter}Let $n = 3(l + 1)$. 
\begin{enumerate}
\item\label{it:stoch}The probability distribution $q$ on $E\times [3]$ defined $($for $i\in [3])$ by 
\begin{align*}\text{ for }j \in [l], \quad q_{u_j u_{j-1}, i} = q_{v_j v_{j -1}, i} &= \frac{1}{6n} \left(2 + (-2)^{- (l - j)}\right) \\
 \text{ for } j \in [l-2], \quad q_{w_j w_{j -1}, i}  &= \frac{1}{6n} \left(2 + (-2)^{- (l - j - 1)}\right) \\
q_{w_{l-1} w_{l -2}, i} = q_{\bar{w}_{l-1} \bar{w}_{l -2}, i} = \frac{1}{2} q_{w_{l} w_{l -1}, i} = \frac{1}{2} q_{\bar{w}_{l} \bar{w}_{l -1}, i} &= \frac{1}{6n} \left(\frac{3}{2}\right)
\end{align*}
on the rows of $M$ has expectation $\frac{1}{n}\id_{3n}$. That is, if the rows of $M$ are scaled by the values of $q$, each column sums to $1/n$. Note that the entries of $q$ are $\Theta(\frac{1}{n})$. Ignoring the index $i$ in $q_{uv,i}$ allows us to view $q$ as a labeling of the edges of the graph $D_l$; see \cref{fig:column-sum,fig:graph}. 
\item\label{it:kernel} $\ker M = \spn (\Omega_0')^\perp$ is spanned by the 2 dimensional space $S \subseteq \RR^{W \times [3]}$ given by
$$ S = \{s: s(v, 1) = \alpha, s(v, 2) = \beta, s(v, 3) = \gamma \text{ for all } v \in W,\; \alpha + \beta + \gamma = 0 \}$$
 and the function $f \in \RR^{W \times [3]}$ which for all $i \in [3]$ assigns
 \begin{align}f (u_j, i) = f(v_j, i) = f(w_j, i) &= (-2)^{-j} \text{ for } j \in [l] \cup \{0\}\nonumber \\
 \textrm{ and } f(\bar{w}_{l-k}, i)  &=  f(w_{l-k}, i) \text{ for } k \in \{0,1\}.\label{eq:kernel} 
 \end{align}
 Note that $f \in (\id_n^\perp)^3 \subseteq S^\perp$. Thus we have the orthogonal decomposition $\spn(\Omega_0')^\perp = S \oplus \spn{f}.$
\item\label{it:singular} Apart from the three zero singular values, all singular values of $M$ are $\Omega(1/n)$.

\end{enumerate}

\end{lem}


Given the lemma, let us prove that the diameter bound holds according to the proof outline at the beginning of the section.

\begin{proof}[Proof of \cref{thm:diameter}] We first show the claim for $n$ of the form $n = 3(l + 1)$; the bound follows for $3(l+1) < n < 3(l + 2)$ by applying \cref{prp:pad-diameter} with $t = 3(l +1)$, using that the array we construct has capacity $1/2$ and $t/n \geq 2/3$. 

We now show the diameter lower bound for $n = 3(l+1)$. It is enough to exhibit a constant $C>0$, and a probability distribution $p$ on $\Omega_{n,3}' = \{e_i: i \in [n]\}^3$ such that for for all $N \geq C n^2 \log n$ and all $x, y, z \in \id_n^\perp$, 
 $$\sum_{\omega \in \Omega_{n,d}'} p_\omega e^{\omega \cdot (x,y,z)} \leq  e^{-N} + \inf_{x', y', z' \in \id_n^\perp} \sum_{\omega \in \Omega_{n,d}'}\; p_\omega e^{\omega \cdot (x,y,z)}$$
 only if $\|(x,y,z)\|_2 = \Omega( 2^{n/3} N)$. Note that the space $(\id_n^\perp)^3$ over which we are infimizing is a subspace of $S^\perp$ where $S$ is as in \cref{lem:diameter}, and that $\Omega_{n,d}' \subseteq S^\perp$. The proof will follow the outline in \cref{subsec:outline}; namely, we will consider a subset $\Omega_0' \subseteq \Omega_{n,d}'$ and an element $\omega' \in \Omega_{n,d}'$ very close to, but outside of, $\aff(\Omega_0')$.

Consider the set $\Omega_0' \subseteq \Omega_{n,d}'$ of rows of $M$ in \cref{lem:diameter} and the probability distribution $q$ on $\Omega_0'$ from \cref{lem:diameter}. Let $\omega' = (e_{u_l}, e_{v_l}, e_{w_l})$ for the vertices $u_l, v_l, w_l \in D_l$.
Let $\Omega = \Omega_0' \cup \{\omega'\}$, and define the probability distribution $p$ on $\Omega$ by $p_{\omega'} = \frac{1}{2}$ and $p_{\omega} = \frac{1}{2} q_{\omega}$ for $\omega \in \Omega_0'$. Recall from \cref{lem:diameter} the orthogonal decomposition $\spn(\Omega_0')^\perp = S \oplus \spn{f}$. As $\RR^{3n} = \spn(\Omega_0') \oplus \spn(\Omega_0')^\perp$, we have the orthogonal decomposition $S^\perp = \spn(\Omega_0') \oplus \spn f$. Observe that $\omega' \not\in \spn(\Omega_0')$, because by \cref{lem:diameter} we have $\spn (\Omega_0')^\perp = \ker M =  S + \spn f$ and clearly $f \cdot \omega' \neq 0$.

By \cref{it:stoch} of \cref{lem:diameter} we have $\sum_{\omega \in \Omega_0'} q_\omega \omega = \frac{1}{n}(\id_n, \id_n, \id_n)$ and thus $\capa(q) = 1$. Therefore, $\omega' \not\in \spn(\Omega_0')$ implies that the infimum is $1/2$ for this choice of $\Omega$ and $p$. We claim that the infimum can only be approximately attained by $h \in (\id_n^\perp)^3$ with a very large component in the one-dimensional space $\spn f = \spn(\Omega_0')^\perp \cap (\id_n^\perp)^3$. As in the proof outline, we must bound the components in $\spn(\Omega_0')$ of the approximate minimizer $h$. For $h \in (\id_n^\perp)^3$ write $h = h_0 + a f$ and $\omega' = \omega_0 + bf$ where $h_0, \omega_0 \in \spn \Omega_0'$. Note that $|b| = \frac{|f \cdot \omega'|}{\|f\|^2} = O(2^{-l}) =  O(2^{-n/3})$ and that $h_0 \in (\id_n^\perp)^3$, because $h$ and $f$ are.
Suppose 
$$ \sum_{\omega \in \Omega} p_{\omega} e^{\omega \cdot h} \leq \frac{1}{2}e^{-N} + \frac{1}{2}.$$
Equivalently, 
\begin{align}
 \sum_{\omega \in \Omega_0'} q_{\omega} e^{\omega \cdot h_0} +  e^{h_0\cdot \omega_0 + a b \norm{f}^2} \leq e^{-N} + 1.\label{eq:almost}
\end{align}
Suppose $\|h_0\|$ is bounded by $L$. If $e^{h_0\cdot \omega_0 + ab \norm{f}^2 } \leq e^{-N}$, then $|a b| =  \Omega(N - L)$. In particular, $\|h\| \geq \|af\| =\vert a b \vert \|f\| / \vert b \vert= \Omega ((N - L)2^{n/3})$ because of the previous bounds on $|ab|, |b|,$ and the fact that $\|f\| = \Theta(1)$. It remains to prove a bound $L$ for $\|h_0\|$. We will do this by showing that if $\|h_0\|$ were too large, then the first term of the left-hand side of \cref{eq:almost} would be too large. This amounts to $\frac{1}{n}\id_{3n}$ being in the relative interior of $\conv(\Omega_0')$, but will be proved using lower bounds on the singular values of $M$.

Let $\alpha$ denote the least nonzero singular value of $M$; by \cref{it:singular} of \cref{lem:diameter} $\alpha = \Omega( 1/n)$. As $h_0 \in \spn(\Omega_0') =\operatorname{rowspan}(M)$, we have $\|M h_0\| \geq \alpha \|h_0\|$ by the singular value bound. We claim that there is some $\omega \in \Omega_0'$ satisfying $\omega \cdot h_0 = \Omega(\alpha \|h_0\|/n)$. To prove this, first note that the $\sum_{\omega \in \Omega_0'} q_{\omega} \omega \cdot h_0 = \frac{1}{n}(\id_n, \id_n, \id_n) \cdot h_0 = 0$ because $h_0 \in  (\id_n^\perp)^3$. Moreover, by \cref{lem:diameter} we have $q_{\omega} = \Theta(1/n)$. The claim follows from \cref{lem:simple} below applied to the sequence $(\omega \cdot h_0: \omega \in \Omega_0').$


Because $q_\omega=\Theta(\frac{1}{n})$, we must have that $\omega \cdot h_0 = O(\log n)$ for all $w \in \Omega_0'$. Else, the contribution from the term $q_\omega e^{\omega \cdot h_0}$ alone is larger than $1$, in which case $x$ cannot be an $e^{- N}$-approximate minimizer. Finally, $\|h_0\| = O(n(\log n)/\alpha) = O(n^{2}\log n),$ and so we may take $L = O(n^{2}\log n)$ and $N \geq 2L$. \end{proof}

In the above proof, we used the following simple lemma.
\begin{lem}\label{lem:simple} Let $0 < \beta < \gamma$. Suppose $z \in \RR^m$ is such that $\sum_{i = 1}^m q_i z_i = 0$ for $q_i \in (\beta/m, \gamma/m)$. Then there exists $i \in [m]$ such that $z_i \geq \frac{\beta }{2\gamma m} \|z\|_2$.
\end{lem}
\begin{proof}
Because $\sum q_i z_i = 0$, 
$$ \sum_{i: z_i < 0} q_i |z_i| = \sum_{i:z_i \geq 0} q_i z_i,$$
and 
$$ \sum_{i: z_i < 0} q_i |z_i| +  \sum_{i:z_i \geq 0} q_i z_i \geq (\beta/m) \|z\|_1 \geq (\beta/m)\|z\|_2.$$
Thus $\sum_{i: z_i \geq 0} q_i |z_i| \geq \frac{\beta}{2m}\|z\|_2,$ so there is some $i$ such that $q_i z_i > \frac{1}{m} \frac{\beta}{2m}\|z\|_2$. Thus $z_i > \frac{\beta}{2\gamma m}\|z\|_2$.
\end{proof}
To show that our diameter lower bound holds for all values of $n$, we need the following proposition, which is proved in \cref{sec:rounding}. The idea is to prove diameter bounds for larger arrays from diameter bounds for smaller ones by embedding the smaller array in a ``corner'' of the larger array.

\begin{prp} \label{prp:pad-diameter}
Suppose $1 \leq t \leq n$. Let $p$ be a $d$-dimensional array in $(\RR_{\geq 0}^t)^{\ot d}$ with unit sum; in particular $\capa(p) \leq 1$. Let $q$ be the $d$-dimensional array in $(\RR_{\geq 0}^n)^{\ot d}$ array such that $q_{i_1, \dots, i_d} = \frac{t}{n} p_{i_1, \dots, i_d}$ for $i_1, \dots, i_d \in [t]$, $q_{iii} = 1/n$ for $t+1 \leq i \leq n$, and $q_{i_1, \dots, i_d} = 0$ otherwise. For $\eps \leq 1 - \capa(p)$, 
$$D_{f_q}(\eps) \geq D_{f_p}\left(\frac{(1 - \capa(p))\eps}{1 - \capa(p)^{t/n}}\right).$$ 
In particular, the norm of any $\eps$-approximate minimizer of $f_q$ is at least the norm of some $\big( \frac{1 - \capa(p)}{1 - \capa(p)^{t/n}} \big) \eps$-approximate minimizer of $f_p$. 
\end{prp}



As a corollary of the proof of \cref{thm:diameter}, we have a bound on the \emph{facet gap} of \cite{burgisser2020interior}. The facet gap of a finite set $\Omega$ is defined to be the least distance of an element of $\Omega$ to the affine hull of a facet of $\conv(\Omega)$. We have shown that the distance between $\aff(\Omega_0')$ and $\omega'$ is $O(2^{-l})$, or $O(2^{- n/3})$. 
\begin{cor}[Facet gap of array scaling] \label{cor:facet-fap}
There is a subset $\Omega_1 \subseteq \Omega_{n,3}$ with facet gap $O(2^{- n/3})$.
\end{cor}

Analogously to what is done for the margin in Proposition C.1 of the original publication, we may also embed this array inside a larger array to obtain a diameter bound for $d \geq 3$. For $d \geq 3$, take $q(i,j,k,l,l,\dots, l) = \frac{1}{n} p_{ijk}$ for all $i,j,k,l \in [n]$. Then for $(x_1, \dots, x_d) \in (\id_n^\perp)^d$ we have 
$$f_q(x_1, \dots, x_d) = \frac{1}{n} f_p(x_1,x_2,x_3)  \sum_{l = 1}^n e^{\sum_{j = 4}^d (x_j)_l}.$$
For fixed $x_1, x_2, x_3$, by Jensen's inequality $f_q$ is minimized when $x_j = 0_n$ for $j \geq 4$ and takes value $f_p(x_1, x_2, x_3)$, and thus $f_q$ has the same diameter bound as $f_p$. 


\begin{cor}[Diameter bound for $d \geq 3$]\label{cor:diameter-d}
There is an absolute constant $C > 0$ such that the following holds. For all $d \geq 3$, there is a family of arrays $q \in (\RR_{\geq 0}^n)^{\ot d}$ with $O(n^2)$ nonzero entries, each of bit-complexity $O(n)$, that satisfies the following property. For all $0 <\eps \leq  \exp(- C n^2 \log n)$ and $x \in \RR^{dn}$, if
	$$f_q (x) \leq \capa(p) + \eps$$
	then $\norm{x}_2 = \Omega\left(2^{n/3}\log(1/\eps)\right).$

\end{cor}



\subsection{Proof of the properties of the construction}

 We now prove \cref{lem:diameter}.
 \begin{proof}[Proof of \cref{lem:diameter}] It is first helpful to change basis on each copy of $\RR^3$ so that the $A$ blocks are diagonalized. Let $U \in \Mat(3)$ be an orthogonal matrix such that 
$$U^\dagger A U =   \begin{bmatrix} 2 & 0 & 0 \\ 0 & -1 & 0 \\ 0 & 0 & -1 \end{bmatrix}.$$
This is possible because $2, -1, -1$ are the eigenvalues of the symmetric matrix $A$. In particular, the first column of $U$ is $(1,1,1)/\sqrt{3}$, and the second two columns span the space of vectors with sum zero. Then $M' = (U^{\oplus n})^\dagger M U^{\oplus n}$ is of the form $P \oplus L \oplus L$ where $P_{e,v} = M'_{(e, 1), (v, 1)}$ for $e, v \in E \times V$ and $L_{e, v} = M'_{(e, 2), (v, 2)}$.  Note that $L$ is the edge-vertex incidence matrix of the directed graph $D_l$, the row corresponding to the edge $(u,v)$ of $D_l$ has a $-1$ in the column indexed by the vertex $u$ and a $+1$ in the column indexed by $v$. Moreover, $P$ is the matrix obtained from $L$ by replacing every $-1$ entry by a $2$.

To prove \cref{it:kernel}, observe that $\ker M$ is $(U^{\oplus n}) \ker M' = (U^{\oplus n}) \ker P \oplus \ker L \oplus \ker L$. Because $D_l$ is connected, $\ker L = \spn \id_n$. As the second two columns of $U$ span the subspace of $\RR^3$ of vectors with sum $0$, the two-dimensional space $S$ is given by $(U^{\oplus n}) 0 \oplus \spn \id_n \oplus \spn \id_n = (U^{\oplus n}) 0 \oplus \ker L \oplus \ker L$. We next reason for $\ker P$, the other summand of the orthogonal decomposition of $\ker M'$. The graph $D_l$ is a connected tree, so $\ker P$ is one dimensional. This is because every choice of $g(w_0) \in \RR$ determines a unique function $g:V \to \RR$ in $\ker P$. We claim that the function $g(v) = f(v,1)$ for $f$ as in \cref{eq:kernel} is in $\ker P$, and hence spans it. To check this, one must check that for every edge $(v,w) \in E$ we have $2g(v) + g(w) = 0$. It is instructive to look at \cref{fig:graph}. Observe that this property holds for the edges $u_{k, k-1}$ if the sequence $g(u_k)$ obeys the recurrence relation $g(u_{k-1}) = -2g(u_k)$ for $k \in [l]$, which is indeed true by the definition of $f$. Checking the condition for $v$ and $w$ is similar. As the first column of $U$ is proportional to $\id_3$, $(U^{\oplus n}) \ker P \oplus 0 \oplus 0$ is spanned by the function $f$. This proves \cref{it:kernel}.

To show \cref{it:singular}, it is enough to argue that the singular values of $P, L, L$ obey the desired bound. For $L$ this follows straightforwardly from the fact that $L$ is an incidence matrix of a connected, directed tree and so is totally unimodular with linearly independent rows. The singular value bound follows by \cref{lem:unimod_sing}. Rather than arguing spectrally for $P$, we make an ad-hoc argument using the structure of $D_l$.
We first show that $\|x^t P\|_\infty = \Omega( \|x\|_\infty)$ for all $x \in \RR^{n-1}$, which suffices because $\|x^t P\|_2 \geq  \|x^tP\|_\infty$ and $\|x\|_\infty \geq \frac{1}{\sqrt{n}}\|x\|_2$.

Let $x \in \RR^{n-1} \setminus \lbrace 0 \rbrace$ and $e$ be an edge in $D_l$ such that $\vert x(e) \vert  = \| x \|_\infty$. If $e = u_iu_{i-1}$ for $i \in [l]$, then $|x^t P (u_{i})| \geq \| x \|_\infty$ because either $i < l$ in which case
	\begin{align*}
	\vert x^t P (u_{i }) \vert = \vert 2 x(u_{i}u_{i - 1}) + x(u_{i+1}u_{i}) \vert \geq 2\|x\|_\infty - \vert x(u_{i+1}u_{i}) \vert \geq \|x\|_\infty
	\end{align*}
or $i = l$ and so $\vert x^t P (u_{i}) \vert = \vert 2 x(e) \vert = 2 \|x\|_\infty$. The same argument applies to all other edges except $e = w_{l-2} w_{l-3}$. In the latter case we are done if $x^t P (w_{l-2}) \geq 1/3 \| x \|_\infty$. Otherwise we necessarily have $\vert x(w_{l-1} w_{l-2}) \vert + \vert x(\bar{w}_{l-1} w_{l-2}) \vert \geq 5/3 \|x\|_\infty$, since $x^t P(w_{l-2}) = 2x(e) + x(w_{l-1} w_{l-2}) + x(\bar{w}_{l-1} w_{l-2})$. It follows that $\vert x(w_{l-1} w_{l-2}) \vert \geq 5/3 \|x\|_\infty - \vert x(\bar{w}_{l-1} w_{l-2}) \vert \geq 5/3 \|x\|_\infty - \|x\|_\infty \geq 2/3 \|x\|_\infty$. As $\vert x^t P(w_{l-1}) \vert = 2  x(w_{l-1} w_{l-2})+ x(w_l w_{l-1}) $, we have 
	\begin{align*}
	\vert x^t P(w_{l-1}) \vert \geq 2 \vert x(w_{l-1} w_{l-2}) \vert - \vert x(w_l w_{l-1}) \vert \geq \frac{4}{3} \|x\|_\infty - \|x\|_\infty \geq \frac{1}{3} \|x\|_\infty.
	\end{align*}
In any case, there is some value of $x^t P$ with absolute value greater or equal $1/3 \|x\|_\infty$. 


Finally, for \cref{it:stoch} we note that the probability distribution $q$ on the rows of $M$ has expectation equal to the all $1/n$ function if and only if the probability distribution $q'$ defined by $q'_e = 3q_{e, 1}$ on the rows of $P$ has expectation equal to the all $3/n$ function on the vertices of $D_l$. Recall that $P$ is obtained from the edge-vertex incidence matrix of $D_l$ by replacing every $-1$ with a $2$. Thus the expectation of the rows under $q'$ at a vertex $v$ is $\sum_{w: (w,v) \in D_l} q'_{(w,v)} + \sum_{w: (v,w) \in D_l} 2q'_{(v,w)}$; see \cref{fig:column-sum}. We now check that this is equal to $3/n$ for each vertex of $D_l$; it is helpful to look at \cref{fig:graph}. The leaves $u_l, v_l, w_l,$ and $\overline{w}_l$ all have outdegree one and indegree zero, and $q'$ takes the value $3 \cdot 3/6n = 3/2 n$ on the outgoing edges. The expectation under $q'$ thus takes value $3/n$ on these vertices. On vertices of indegree one and outdegree one, $q'$ takes the value $\frac{1}{2n} \left(2 + (-2)^{- k}\right)$ on the incoming edge and $\frac{1}{2n} \left(2 + (-2)^{- (k+1)}\right)$ on the outgoing edge. Thus the expectation takes the value $\frac{1}{2n} \left(2 + (-2)^{- k}\right) + \frac{1}{2n} \left(4 - (-2)^{- k}\right) = 3/n.$ The remaining vertices to check, those of total degree three, are $r$ and $w_{l-2}$. For $r$, which has only incoming edges, the expectation under $q'$ is $2\cdot \frac{1}{2n} \left(2 + (-2)^{- (l-1)}\right) + \frac{1}{2n} \left(2 + (-2)^{- (l-2)}\right),$ which is again $3/n$. For $w$ the expectation is $2\cdot \frac{1}{2n} \left(2 -\frac{1}{2}\right) + 2\cdot \frac{1}{2n} \left(2 -\frac{1}{2}\right) = 3/n.$ This completes the proof. \end{proof}



\begin{lem}\label{lem:unimod_sing}
If $A$ is an $n \times k$ totally unimodular matrix with linearly independent columns, then the eigenvalues of $A^T A$ are all at least $1/n^2$. 

\end{lem}

\begin{proof} First note that $k \leq n$ by the linear independence of the columns of $A$. The least eigenvalue of $A^T A$ is $\min_{x \in \RR^k\setminus \{0\}} (x^T A^T A x)/\|x\|^2 = \min_{x \in \RR^k\setminus \{0\}} \|Ax\|^2/\|x\|^2$, so it suffices to show that for all $x \in \RR^k$, $A x$ has norm at least $\|x\|/n$. Indeed, if $Ax = y$, then there is some invertible $k\times k$ submatrix $A'$ of $A$ and $k \times 1$ submatrix $y'$ of $y$ such that $A'x = y'$. By Cramer's rule and unimodularity of $A'$ we have that, for $i \in [k]$,
	$$ x_i = \dfrac{\det(B_i)}{\det(A')} = \pm \det(B_i)$$
	where $B_i$ is simply the matrix that one obtains by replacing the $i^{th}$ column of $A'$ with the
	vector $y'$. By performing the Laplace expansion with respect to the $i^{th}$ column, and by unimodularity of the
	minors, we have that $x_i \leq \|y\|_1$, and so $\|x\|_2 \leq \sqrt{k} \|y\|_1 \leq n \|Ax\|_2$ (using $k \leq n$).
\end{proof}


%





















































































































































































































































































































































































































































































































































































































































































































































































































































































































































\appendix\setcounter{section}{4}\section{Padding and rounding for diameter bounds}\label{sec:rounding}

We begin with the proof of \cref{prp:pad-diameter}. We prove it only for $d = 3$, but the proof goes through mutatis mutandis for all $d \geq 1$.

\begin{proof}[Proof of \cref{prp:pad-diameter}]
Recall that $q$ is the $n\times n \times n$ array such that $q_{ijk} = \frac{t}{n} p_{ijk}$ for $i,j,k \in [t]$, $q_{iii} = 1/n$ for $t+1 \leq i \leq n$, and $q_{ijk} = 0$ otherwise. We may split the inputs $x,y,z \in \id_n^\perp$ into 
\begin{align*}x & =\left(x' + \alpha_1 \id_t ,x'' - \frac{t}{n-t} \alpha_1 \id_{n-t}\right),\\
 y&=\left(y' + \alpha_2 \id_t, y'' - \frac{t}{n-t} \alpha_2 \id_{n-t}\right),\\
 z& = \left(z'+ \alpha_3 \id_t ,z''- \frac{t}{n-t} \alpha_3 \id_{n-t}\right)
\end{align*} where $x',y',z' \in \RR^{t}, x'', y'', z'' \in \RR^{n-t}$ each sum to zero; write $w = (x',y',z')$. As $\| (x,y,z)\|_2 \geq \|w\|_2$, it is enough to prove that $\|w\|_2$ is large for any approximate minimizer. By optimizing over $\alpha_i$ and $x'', y'', z''$ for fixed $w$, one computes that the optimum value for $f_q$ for any fixed $w$ is $f_p(w)^{t/n}$. To see this, write
$$ f_q(x,y,z) = \frac{t e^{\alpha_1 + \alpha_2 + \alpha_3}}{n} f_p(w) + \frac{e^{ -\frac{t}{n - t} (\alpha_1 + \alpha_2 + \alpha_3)} }{n} \sum_{i = t+1}^n e^{x''_i + y''_i + z''_i}.$$
First note that for fixed $\alpha_i$'s, the second term is minimized at $x'' = y'' = z'' = 0$ by Jensen's inequality. Furthermore, the value only depends on $\alpha:= \alpha_1 + \alpha_2 + \alpha_3$. With $x'',y'',z'' = 0$, we have
$$ f_q(x,y,z) =g(w, \alpha):= \frac{t e^{\alpha}}{n} f_p(w) + \frac{(n - t) }{n} e^{- \frac{t}{n - t} \alpha}.$$
Taking the derivative in $\alpha$, we see that this is minimized when $ f_p(w) e^{\alpha} = e^{- \frac{t}{n - t} \alpha},$ or $e^{\alpha} = f_p(w)^{-1/(1 + \frac{t}{n-t})} = f_p(w)^{-\frac{n-t}{n}}.$ Plugging this value in proves that the optimum is $f_p(w)^{t/n}$.
By concavity of $x^{t/n},$ provided $f_p(w) \leq 1$ we have 
$$f_p(w)^{t/n} - \capa(p)^{t/n} \geq \frac{1 - \capa(p)^{t/n}}{1 - \capa(p)} (f_p(w) - \capa(p)).$$
The first factor in the second term is the slope of the line from $(\capa(p), \capa(p)^{t/n})$ to $(1,1)$. Thus for any $\eps \leq 1 - \capa(p)$, any $\eps$-approximate minimizer for $f_q$ has norm at least that of some $\big( \frac{1 - \capa(p)}{1 - \capa(p)^{t/n}} \big)\eps$-approximate minimizer for $f_p$.
\end{proof}









































